\section{问题四的建立与求解}

本题针对女性胎儿样本，由于缺少Y染色体标记，无法沿用前述基于Y染色体浓度的分析路径。需要利用NIPT检测的多维汇总统计量，建立染色体非整倍体异常（13-三体、18-三体、21-三体及性染色体异常）的检测模型。数据集包含604例女性胎儿样本，其中正常537例、异常67例，异常率约为11.1\%。

\subsection{问题转化与特征选择}

女性胎儿不携带Y染色体，前三个问题的Y染色体浓度建模思路不再适用。问题四的核心任务转化为，在多维NIPT汇总指标空间中将异常样本与正常样本区分开来。

候选特征涵盖各目标染色体的Z值、BMI、X染色体浓度、GC含量、读段计数及比对率等。经初步相关性分析与临床意义筛选，确定以下6个核心特征构成输入空间。

\begin{itemize}
  \item 13号染色体Z值（$Z_{13}$）、18号染色体Z值（$Z_{18}$）、21号染色体Z值（$Z_{21}$）
  \item X染色体Z值（$Z_X$）
  \item 孕妇BMI
  \item X染色体浓度
\end{itemize}

该特征集覆盖了三种常见常染色体三体及性染色体异常的直接度量指标，同时纳入BMI和X染色体浓度作为辅助变量，以捕捉母体特征对检测结果的潜在调节作用。

\subsection{多元马氏距离异常检测模型}

\subsubsection{模型建立}

本文建立\textbf{多元马氏距离异常检测模型（记为M1）}作为主方法。其基本思路为，以正常样本的多维分布为参照，计算每个待检样本偏离正常群体中心的马氏距离，距离越大则异常可能性越高。

设正常样本集为$\{\mathbf{x}_k\}_{k=1}^{n_0}$，其中$\mathbf{x}_k \in \mathbb{R}^6$为6维特征向量。首先估计正常群体的均值向量$\boldsymbol{\mu}$和协方差矩阵$\boldsymbol{\Sigma}$
\begin{equation}
  \boldsymbol{\mu} = \frac{1}{n_0}\sum_{k=1}^{n_0}\mathbf{x}_k, \qquad
  \boldsymbol{\Sigma} = \frac{1}{n_0-1}\sum_{k=1}^{n_0}(\mathbf{x}_k - \boldsymbol{\mu})(\mathbf{x}_k - \boldsymbol{\mu})^T
\end{equation}

对任意待检样本$\mathbf{x}$，其马氏距离定义为
\begin{equation}\label{eq:mahal}
  d_M(\mathbf{x}) = \sqrt{(\mathbf{x} - \boldsymbol{\mu})^T \boldsymbol{\Sigma}^{-1} (\mathbf{x} - \boldsymbol{\mu})}
\end{equation}

分类规则为，若$d_M(\mathbf{x}) \geq \tau$，则标记为异常，否则判定为正常。阈值$\tau$取正常样本马氏距离分布的某一百分位数，通过遍历不同百分位（p50至p99）选取F1分数最优的切点。

\subsubsection{求解步骤}

\textbf{Step 1}（特征标准化）\quad 对604例样本的6维特征分别进行均值中心化和方差归一化，消除量纲差异对距离度量的影响。

\textbf{Step 2}（参数估计）\quad 仅使用537例正常样本估计$\boldsymbol{\mu}$和$\boldsymbol{\Sigma}$。同时采用最小协方差行列式（MinCovDet）方法估计稳健协方差矩阵，作为对比方案。

\textbf{Step 3}（距离计算与阈值搜索）\quad 对全部604例样本计算马氏距离，在正常样本距离的第50至第99百分位区间内搜索最优阈值，以F1分数为选取准则。

\textbf{Step 4}（性能评估）\quad 计算召回率、精确率、F1分数、AUC-ROC和AUC-PR，全面评估模型的异常检测能力。

\subsection{基线方法}

本文建立\textbf{多元Z值联合阈值规则（记为M2）}作为基线方法。M2不考虑特征间的相关结构，直接对三条常染色体的Z值取平方和
\begin{equation}\label{eq:zsquare}
  S = Z_{13}^2 + Z_{18}^2 + Z_{21}^2
\end{equation}

分类规则为，若$S \geq \tau$，则标记为异常。阈值$\tau$同样在正常样本$S$值分布的百分位区间内搜索确定。

M2与M1的本质区别在于，M2将协方差矩阵退化为单位阵，仅利用各Z值的边际信息，忽略了特征间的联合分布结构。

\subsection{求解结果}

\subsubsection{M1模型结果}

在经验协方差估计下，M1在正常样本马氏距离第85百分位处取得最优F1分数，对应阈值$\tau=2.997$。表~\ref{tab:q4-m1}列出该阈值下的混淆矩阵与主要评价指标。

\begin{table}[H]
  \centering
  \caption{M1（多元马氏距离异常检测）在最优阈值下的检测结果}\label{tab:q4-m1}
  \begin{tabular}{lcccccc}
    \toprule
    阈值（百分位） & TP & FP & FN & TN & 召回率 & 精确率 \\
    \midrule
    p85（$\tau=2.997$） & 18 & 81 & 49 & 456 & 26.87\% & 18.18\% \\
    \bottomrule
  \end{tabular}
\end{table}

M1的F1分数为0.217，AUC-ROC为0.574，AUC-PR为0.163。采用MinCovDet稳健协方差估计后，性能与经验协方差方案接近，未见显著提升。将特征集扩展至15维（加入GC含量、读段计数、比对率等）后仅获得边际改善；仅保留$Z_{13}$、$Z_{18}$、$Z_{21}$三个特征的精简方案则略逊于6维核心特征集。

图~\ref{fig:q4-roc}展示了不同特征集与协方差估计方式下M1的ROC曲线对比，图~\ref{fig:q4-dist}展示了正常与异常样本的马氏距离分布。

\begin{figure}[H]
  \centering
  \includegraphics[width=0.85\textwidth]{Q4/experiments/round1/figures/q4_m1_roc_comparison.png}
  \caption{M1在不同配置下的ROC曲线对比}\label{fig:q4-roc}
\end{figure}

\begin{figure}[H]
  \centering
  \includegraphics[width=0.85\textwidth]{Q4/experiments/round1/figures/q4_m1_distance_distribution.png}
  \caption{正常样本与异常样本的马氏距离分布}\label{fig:q4-dist}
\end{figure}

\subsubsection{M2基线结果}

M2在正常样本$S$值第50百分位处取得最优性能，召回率为47.8\%，精确率为10.6\%，F1分数为0.174。AUC-ROC仅为0.492，接近随机分类水平。

\subsubsection{M1与M2对比}

表~\ref{tab:q4-compare}汇总了两种方法在关键指标上的对比。

\begin{table}[H]
  \centering
  \caption{M1与M2检测性能对比}\label{tab:q4-compare}
  \begin{tabular}{lccccc}
    \toprule
    模型 & AUC-ROC & AUC-PR & F1 & 召回率 & 精确率 \\
    \midrule
    M1（马氏距离） & 0.574 & 0.163 & 0.217 & 26.87\% & 18.18\% \\
    M2（Z值平方和） & 0.492 & --- & 0.174 & 47.80\% & 10.60\% \\
    \bottomrule
  \end{tabular}
\end{table}

M1的AUC-ROC较M2高出0.082，在精确率-召回率权衡上也更优。M2虽然召回率较高，但精确率仅为10.6\%，大量正常样本被误判为异常，实际可用性不足。M1通过协方差矩阵捕获特征间的联合分布结构，在区分能力上获得了可测量的增益。

\subsection{结果分析}

两种方法的AUC-ROC均处于0.5至0.6区间，整体检测能力有限。这一结果反映了问题本身的内在困难，异常样本与正常样本在NIPT汇总统计量空间中的特征分布高度重叠，多数染色体非整倍体异常仅引起Z值的微小偏移，难以在群体层面形成清晰的分离边界。

M1相对于M2的优势来源于马氏距离对特征间相关结构的利用。当多个Z值同时出现轻微偏移时，马氏距离能够识别这种联合偏离模式，而M2的平方和指标对各维度等权处理，丢失了协方差信息。按异常类型分析，不同三体的检出率存在显著差异，21-三体因Z值偏移幅度相对较大而更容易被捕获，13-三体和性染色体异常的检出难度则更高。

从临床应用角度评估，M1提供了优于随机分类的筛查能力，可作为初步筛选工具缩小需要进一步确认的样本范围。在实际产前诊断流程中，NIPT结果呈阳性的孕妇仍需接受羊膜穿刺或绒毛膜绒毛取样等确证性检查。本模型的定位是在确证检查之前提供风险分层，降低不必要的侵入性检查比例，而非替代金标准诊断。检测能力的根本瓶颈在于NIPT汇总统计量对微小染色体剂量变化的灵敏度不足，提升检测精度需要引入更细粒度的测序特征或结合多平台数据。
